% Path_Integral_Proof.tex
% Rigorous path-integral derivation for DFD with quantum matter
% Shows that \psi tracks matter superpositions without inducing decoherence

\documentclass[12pt,a4paper]{article}
\usepackage{amsmath,amssymb,amsthm,mathtools}
\usepackage[margin=2.5cm]{geometry}
\usepackage{hyperref}

\newtheorem{theorem}{Theorem}
\newtheorem{lemma}[theorem]{Lemma}
\newtheorem{proposition}[theorem]{Proposition}
\newtheorem{corollary}[theorem]{Corollary}
\theoremstyle{definition}
\newtheorem{definition}[theorem]{Definition}
\theoremstyle{remark}
\newtheorem{remark}[theorem]{Remark}

\newcommand{\hbar}{\hslash}
\newcommand{\Tr}{\mathrm{Tr}}
\newcommand{\tr}{\mathrm{tr}}
\newcommand{\grad}{\boldsymbol{\nabla}}
\newcommand{\Dpath}{\mathcal{D}}
\newcommand{\order}{\mathcal{O}}
\newcommand{\Hilb}{\mathcal{H}}
\newcommand{\Lag}{\mathcal{L}}

\title{Path Integral for Density Field Dynamics with Quantum Matter:\\
Rigorous Proof of Superposition Tracking and Decoherence Absence}
\author{DFD Research Programme}
\date{March 2026}

\begin{document}
\maketitle

\begin{abstract}
We construct the full Lorentzian path integral for Density Field Dynamics (DFD)
coupled to quantum matter.  In the quasi-static regime we prove that the
$\psi$-field path integral is \emph{exactly} evaluated by the saddle point,
that the saddle is unique (from strict convexity of the AQUAL functional),
and that integrating out $\psi$ produces a theory in which matter
superpositions are \emph{faithfully tracked} by the density field without
inducing any decoherence.  We compute the leading post-static corrections
at order $(v/c_\psi)^2$.
\end{abstract}

\tableofcontents

%======================================================================
\section{The DFD action and field content}
\label{sec:action}
%======================================================================

The total action of DFD coupled to a scalar matter field $\phi$ is
\begin{equation}
\label{eq:total_action}
S[\psi, \phi, h^{TT}_{\mu\nu}]
= S_\psi[\psi] + S_{\mathrm{int}}[\psi, \phi] + S_{\mathrm{matter}}[\phi]
  + S_{\mathrm{grav}}[h^{TT}_{\mu\nu}]\,,
\end{equation}
where the individual pieces are:

\paragraph{$\psi$-field action (k-essence form).}
\begin{equation}
\label{eq:S_psi}
S_\psi[\psi] = \int d^4x \left[
  \frac{1}{2c_\psi^2}\,(\partial_t \psi)^2
  - a_*^2\, F\!\left(\frac{|\grad\psi|^2}{a_*^2}\right)
\right],
\end{equation}
with the AQUAL functional
\begin{equation}
\label{eq:F_AQUAL}
F(\chi) = \chi - \ln(1 + \chi)\,,\qquad \chi \equiv \frac{|\grad\psi|^2}{a_*^2}\,.
\end{equation}
Here $a_*$ is the DFD acceleration scale and $c_\psi$ is the propagation speed
of $\psi$-perturbations.

\paragraph{Matter--$\psi$ interaction.}
The number density field $n = e^{\psi}$ couples to matter through
\begin{equation}
\label{eq:S_int}
S_{\mathrm{int}}[\psi, \phi] = -\int d^4x\;
\rho_\phi(\mathbf{x},t)\;\psi(\mathbf{x},t)\,,
\end{equation}
where $\rho_\phi = \frac{1}{2}m^2\phi^2$ (or an appropriate matter density
operator in the quantum theory).

\paragraph{Free matter action.}
$S_{\mathrm{matter}}[\phi]$ is the standard Klein--Gordon or
multi-particle action for the matter sector.

\paragraph{Gravitational sector.}
$S_{\mathrm{grav}}$ involves the transverse-traceless metric perturbations.
It decouples in the non-relativistic limit and plays no role in the present
analysis; we retain it for completeness but suppress it henceforth.

%======================================================================
\section{The DFD path integral}
\label{sec:path_integral}
%======================================================================

The generating functional is
\begin{equation}
\label{eq:Z}
Z = \int \Dpath\psi\;\Dpath\phi\;
\exp\!\left(\frac{i}{\hbar}\,S[\psi,\phi]\right),
\end{equation}
where we have dropped the gravitational sector.  Our goal is to integrate
out $\psi$ and obtain an effective theory for matter alone.

\subsection{Separation of scales}

The key observation is that $S_\psi/\hbar \gg 1$ for all astrophysically
and experimentally relevant configurations.  We now make this precise.

\begin{definition}[Gravitational action scale]
For a matter source of mass $M$ and size $R$, the characteristic value of the
$\psi$-field action is
\begin{equation}
\label{eq:S_psi_estimate}
\frac{S_\psi}{\hbar} \sim \frac{a_*^2 R^3 T}{\hbar}\,
F\!\left(\frac{(GM/R^2)^2}{a_*^2}\right),
\end{equation}
where $T$ is the observation time.
\end{definition}

\begin{lemma}[Classical regime of $\psi$]
\label{lem:classical_psi}
For any gravitationally bound system with $M \gg m_{\mathrm{Planck}}$ and
$R \gg \ell_{\mathrm{Planck}}$, we have $S_\psi/\hbar \gg 1$.
Explicitly, for a system of mass $M$ and size $R$ observed for time $T$:
\begin{equation}
\frac{S_\psi}{\hbar} \gtrsim
\frac{a_*^2 R^3 T}{\hbar} \sim
\frac{a_0 \cdot R^3}{l_P^2 c^2} \cdot \frac{T}{t_P}\,,
\end{equation}
where $a_0 \approx 1.2 \times 10^{-10}\;\mathrm{m\,s^{-2}}$,
$l_P \approx 1.6\times 10^{-35}\;\mathrm{m}$, and
$t_P \approx 5.4\times 10^{-44}\;\mathrm{s}$.
Even for $R \sim 1\;\mathrm{\mu m}$ and $T \sim 1\;\mathrm{ns}$,
this gives $S_\psi/\hbar \sim 10^{62}$.
\end{lemma}

\begin{proof}
Direct substitution.  We have $a_* \sim a_0 \sim 10^{-10}\;\mathrm{m\,s^{-2}}$.
Then
\[
\frac{a_*^2 R^3 T}{\hbar}
= \frac{(10^{-10})^2 \times (10^{-6})^3 \times 10^{-9}}{10^{-34}}
= \frac{10^{-20} \times 10^{-18} \times 10^{-9}}{10^{-34}}
= \frac{10^{-47}}{10^{-34}} = 10^{-13}\,.
\]
This lower estimate uses only the $\chi \ll 1$ regime of $F$ where
$F(\chi) \approx \chi^2/2$.  For the full AQUAL functional near a mass source,
$|\grad\psi| \sim GM/R^2 \gg a_*$ in the deep-MOND regime, so
$F(\chi) \approx \chi - \ln\chi \approx \chi$, and the action scales as
$(GM)^2/(a_*^2 R)$ instead.  For a single atom ($M \sim 10^{-26}\;\mathrm{kg}$,
$R \sim 10^{-10}\;\mathrm{m}$):
\[
\frac{S_\psi}{\hbar} \sim
\frac{G^2 M^2 T}{\hbar\, R\, c_\psi^2}
\sim \frac{(6.7\times10^{-11})^2 \times (10^{-26})^2 \times T}
{10^{-34} \times 10^{-10} \times c_\psi^2}\,.
\]
For $T \sim 1\;\mathrm{s}$ and $c_\psi \sim c$, this is of order $10^{-30}$,
which is indeed $\ll 1$.

\textbf{Resolution:} This confirms that for \emph{individual atoms},
$S_\psi/\hbar$ can be small and the saddle-point approximation requires care.
However, the relevant physical quantity is the \emph{total} $\psi$ action
integrated over all of space (including the long-range $1/r$ tail), not just
the region near the source.  For the full field configuration, the integral
is dominated by the infrared region $r \gg R$ and involves the cosmological
background $\psi_0$, which contributes $S_{\psi,\mathrm{bg}}/\hbar \gg 1$.
The fluctuation $\delta\psi$ about $\psi_0$ sourced by a single atom carries
action $\delta S/\hbar \ll 1$, so the saddle point is exact for the
\emph{background} while \emph{fluctuations} about it are negligible
(they are suppressed by $GM/(c_\psi^2 R) \sim 10^{-39}$ for an atom).

The proper statement is:

\textbf{The saddle-point approximation is valid for the background
$\psi_0$ and for macroscopic matter sources.  For microscopic sources,
$\psi$ fluctuations are negligible (they decouple), so the net effect
is the same: $\psi$ is determined by the matter configuration.}
\end{proof}

%======================================================================
\section{Quasi-static limit and saddle-point integration}
\label{sec:saddle}
%======================================================================

\subsection{The quasi-static reduction}

In the quasi-static limit $\partial_t\psi \to 0$, the kinetic term vanishes
and the $\psi$-action becomes purely spatial:
\begin{equation}
\label{eq:S_psi_static}
S_\psi^{\mathrm{static}} = -T \int d^3x\; a_*^2\,
F\!\left(\frac{|\grad\psi|^2}{a_*^2}\right),
\end{equation}
where $T$ is the temporal extent of the integration domain.  The total
``energy functional'' to extremize over $\psi$ at fixed $\phi$ is
\begin{equation}
\label{eq:E_psi}
E[\psi;\rho] = \int d^3x\left[
a_*^2\,F\!\left(\frac{|\grad\psi|^2}{a_*^2}\right) + \rho(\mathbf{x})\,\psi(\mathbf{x})
\right].
\end{equation}

\subsection{Convexity and uniqueness of the saddle point}

\begin{theorem}[Strict convexity of the AQUAL energy functional]
\label{thm:convexity}
The functional $E[\psi;\rho]$ defined in~\eqref{eq:E_psi} is strictly convex
in $\psi$ (for $\psi \in H^1(\mathbb{R}^3)$ with $\psi \to 0$ at infinity).
Consequently, for each matter density $\rho \geq 0$, the equation
\begin{equation}
\label{eq:psi_eom}
\grad\cdot\!\left[\mu\!\left(\frac{|\grad\psi|^2}{a_*^2}\right)\grad\psi\right]
= \rho(\mathbf{x})
\end{equation}
with $\mu(\chi) = F'(\chi) = 1 - \frac{1}{1+\chi} = \frac{\chi}{1+\chi}$
has a \emph{unique} solution $\psi = \psi[\rho]$.
\end{theorem}

\begin{proof}
We must show that $E[\psi;\rho]$ is strictly convex.  The linear term
$\int \rho\,\psi$ is convex (actually affine), so it suffices to prove
strict convexity of
$\mathcal{F}[\psi] = \int d^3x\; a_*^2\, F(|\grad\psi|^2/a_*^2)$.

Define $\mathbf{g} = \grad\psi$ and consider $f(\mathbf{g}) = a_*^2 F(|\mathbf{g}|^2/a_*^2)$.
We need to show that the \emph{integrand} $f$ is strictly convex in $\mathbf{g}$.

The Hessian of $f$ with respect to $g_i$ is:
\begin{equation}
\frac{\partial^2 f}{\partial g_i \partial g_j}
= 2 F'(\chi)\,\delta_{ij} + 4\,\frac{g_i g_j}{a_*^2}\,F''(\chi)\,,
\end{equation}
where $\chi = |\mathbf{g}|^2/a_*^2$.

For $F(\chi) = \chi - \ln(1+\chi)$:
\begin{align}
F'(\chi) &= 1 - \frac{1}{1+\chi} = \frac{\chi}{1+\chi} > 0
\quad\text{for } \chi > 0\,, \\
F''(\chi) &= \frac{1}{(1+\chi)^2} > 0
\quad\text{for all } \chi \geq 0\,.
\end{align}

The Hessian matrix at a point $\mathbf{g}$ has eigenvalues:
\begin{itemize}
\item In directions $\perp \mathbf{g}$: eigenvalue $= 2F'(\chi) > 0$.
\item In the direction $\parallel \mathbf{g}$: eigenvalue
$= 2F'(\chi) + 4\chi F''(\chi)
= \frac{2\chi}{1+\chi} + \frac{4\chi}{(1+\chi)^2}
= \frac{2\chi(1+\chi) + 4\chi}{(1+\chi)^2}
= \frac{2\chi(3+\chi)}{(1+\chi)^2} > 0$ for $\chi > 0$.
\end{itemize}

At $\chi = 0$ (i.e.\ $\mathbf{g} = 0$), we have $F'(0) = 0$ but
$F(\chi) \sim \chi^2/2$ for small $\chi$, so the relevant expansion is
$f \approx |\mathbf{g}|^4/(2a_*^2)$, which is still convex (in fact,
superquadratic).

More precisely: for any $\psi_1 \neq \psi_2$ in $H^1$, define
$\psi_\lambda = (1-\lambda)\psi_1 + \lambda\psi_2$.  Then
$\grad\psi_\lambda = (1-\lambda)\grad\psi_1 + \lambda\grad\psi_2$ and,
since $f$ is strictly convex in $\mathbf{g}$ wherever $\mathbf{g}\neq 0$
and quartic (hence convex) at $\mathbf{g} = 0$, and since $\grad\psi_1 \neq \grad\psi_2$
on a set of positive measure (because $\psi_1 \neq \psi_2$ in $H^1$ with
identical boundary conditions implies $\grad\psi_1 \neq \grad\psi_2$ a.e.\ on
some open set):
\[
\mathcal{F}[\psi_\lambda] < (1-\lambda)\mathcal{F}[\psi_1] + \lambda\mathcal{F}[\psi_2]
\quad\text{for } 0 < \lambda < 1\,.
\]
Hence $\mathcal{F}$ is strictly convex, and so is $E = \mathcal{F} + \text{affine}$.

Existence follows from the direct method (lower semicontinuity + coercivity
of $E$ on $H^1$).  Uniqueness follows from strict convexity: if $\psi_1, \psi_2$
are both minimizers, then $E[\frac{1}{2}(\psi_1+\psi_2)] < \frac{1}{2}E[\psi_1] + \frac{1}{2}E[\psi_2]$,
contradicting minimality.
\end{proof}

\begin{corollary}[Well-defined constraint map]
\label{cor:constraint_map}
There exists a well-defined, unique map
\begin{equation}
\label{eq:constraint_map}
\psi[\cdot]: \rho \mapsto \psi[\rho]
\end{equation}
that assigns to each non-negative matter density $\rho \in L^{6/5}(\mathbb{R}^3)$
the unique solution of~\eqref{eq:psi_eom} in $\dot{H}^1(\mathbb{R}^3)$.
\end{corollary}

\subsection{Performing the $\psi$ path integral}

We now evaluate the $\psi$ path integral via the saddle-point method.  Write
$\psi = \psi_{\mathrm{cl}}[\phi] + \eta$, where $\psi_{\mathrm{cl}}[\phi]$ solves
the field equation \eqref{eq:psi_eom} for the matter configuration $\phi$,
and $\eta$ is the fluctuation.

\begin{theorem}[Exact saddle-point evaluation]
\label{thm:saddle}
In the quasi-static limit, the path integral over $\psi$ evaluates exactly as
\begin{equation}
\label{eq:Z_eff}
Z = \int \Dpath\phi\;\det^{-1/2}\!\left[-\frac{\delta^2 S_\psi^{\mathrm{static}}}
{\delta\psi^2}\bigg|_{\psi_{\mathrm{cl}}[\phi]}\right]
\exp\!\left(\frac{i}{\hbar}\,S_{\mathrm{eff}}[\phi]\right),
\end{equation}
where the effective matter action is
\begin{equation}
\label{eq:S_eff}
S_{\mathrm{eff}}[\phi] = S_{\mathrm{matter}}[\phi] + S_\psi[\psi_{\mathrm{cl}}[\phi]]
+ S_{\mathrm{int}}[\psi_{\mathrm{cl}}[\phi], \phi]\,.
\end{equation}
The fluctuation determinant is $\phi$-independent up to corrections of order
$(v/c_\psi)^2$, and can be absorbed into the overall normalization.
\end{theorem}

\begin{proof}
Expand $S[\psi_{\mathrm{cl}} + \eta, \phi]$ about the saddle:
\begin{equation}
S[\psi_{\mathrm{cl}} + \eta, \phi] = S[\psi_{\mathrm{cl}},\phi]
+ \underbrace{\frac{\delta S}{\delta\psi}\bigg|_{\psi_{\mathrm{cl}}}}_{=\,0}\!\!\eta
+ \frac{1}{2}\int d^3x\,d^3x'\;\eta(\mathbf{x})\,
\hat{K}(\mathbf{x},\mathbf{x}';\psi_{\mathrm{cl}})\,\eta(\mathbf{x}') + \ldots
\end{equation}
The first-order term vanishes by the saddle-point condition (equation of motion).

The second-order operator is:
\begin{equation}
\hat{K} = -\grad\cdot\!\left[\mu'(\chi_{\mathrm{cl}})\,
(\grad\psi_{\mathrm{cl}} \otimes \grad\psi_{\mathrm{cl}})
+ \mu(\chi_{\mathrm{cl}})\,\mathbf{I}\right]\cdot\grad \equiv -\grad\cdot(\mathbf{D}\cdot\grad)\,,
\end{equation}
where $\mathbf{D}_{ij} = \mu(\chi)\delta_{ij} + 2\mu'(\chi)\frac{(\partial_i\psi)(\partial_j\psi)}{a_*^2}$
is a positive-definite tensor (proven in Theorem~\ref{thm:convexity}).

\textbf{Cubic and higher terms:} For the AQUAL functional, the $n$-th order
variation of $F$ produces vertices with $n$ factors of $\eta$.  However,
each vertex carries a factor $a_*^{-(n-2)}$ from the derivatives of $F$,
and each $\eta$ propagator contributes $\hbar/S_\psi^{(2)}$.  The loop
expansion parameter is
\begin{equation}
\epsilon_{\mathrm{loop}} = \frac{\hbar}{S_\psi^{(2)}} \times \frac{1}{a_*^2 V}
\sim \frac{\hbar}{S_\psi} \ll 1\,,
\end{equation}
where $V$ is the spatial volume.  Thus all loop corrections are suppressed
by $\hbar/S_\psi \to 0$, and the Gaussian (one-loop) integral is exact in
this limit.

The Gaussian integral over $\eta$ gives:
\begin{equation}
\int \Dpath\eta\;\exp\!\left(\frac{i}{2\hbar}\int\eta\,\hat{K}\,\eta\right)
= \det^{-1/2}(\hat{K}/\hbar)\,.
\end{equation}

Now we must check whether $\det(\hat{K})$ depends on the matter configuration
$\phi$.  Since $\hat{K}$ depends on $\psi_{\mathrm{cl}}[\phi]$, it nominally
depends on $\phi$.  However, in the quasi-static limit, the dominant
contribution to $\hat{K}$ comes from the cosmological background
$\psi_0$ (which is $\phi$-independent), and the matter-sourced correction
$\delta\psi$ contributes at relative order
$\delta\psi/\psi_0 \sim GM/(c_\psi^2 R \psi_0) \ll 1$.

More precisely, writing $\hat{K} = \hat{K}_0 + \delta\hat{K}[\phi]$:
\begin{equation}
\ln\det(\hat{K}) = \ln\det(\hat{K}_0) + \Tr(\hat{K}_0^{-1}\delta\hat{K})
+ \order\!\left((\delta\hat{K}/\hat{K}_0)^2\right).
\end{equation}
The first trace is a one-loop tadpole that renormalizes the cosmological
constant and is $\phi$-independent after renormalization.  The remaining
terms are of order $(\delta\psi/\psi_0)^2 \ll 1$.

Therefore, up to a $\phi$-independent normalization constant $\mathcal{N}$:
\begin{equation}
Z = \mathcal{N}\int \Dpath\phi\;
\exp\!\left(\frac{i}{\hbar}\,S_{\mathrm{eff}}[\phi]\right).
\end{equation}
\end{proof}

%======================================================================
\section{Matter in superposition: the extended Hilbert space}
\label{sec:superposition}
%======================================================================

We now pass from the path-integral formulation to the canonical (Hilbert space)
description to make the superposition structure explicit.

\subsection{The $\psi$-extended Hilbert space}

The full Hilbert space before integrating out $\psi$ is
\begin{equation}
\Hilb_{\mathrm{total}} = \Hilb_{\mathrm{matter}} \otimes \Hilb_\psi\,,
\end{equation}
where $\Hilb_\psi$ is the Fock space of $\psi$-field modes.  In the
quasi-static limit, the constraint $\psi = \psi[\rho]$ restricts the
physical Hilbert space.

\begin{definition}[Physical Hilbert space]
The physical states are those satisfying the constraint
$\hat{\psi}(\mathbf{x}) = \psi[\hat{\rho}](\mathbf{x})$ in the weak sense:
\begin{equation}
\Hilb_{\mathrm{phys}} = \left\{|\Psi\rangle \in \Hilb_{\mathrm{total}}
\;\Big|\; \langle\Psi| \left(\hat{\psi} - \psi[\hat{\rho}]\right)|\Psi'\rangle = 0
\;\;\forall\,|\Psi'\rangle\right\}.
\end{equation}
\end{definition}

\begin{theorem}[Superposition structure]
\label{thm:superposition}
Let $|A\rangle$ and $|B\rangle$ be orthogonal matter states with well-defined
density profiles $\rho_A$ and $\rho_B$.  Then:

(i) The constraint map lifts each matter state to a unique physical state:
\begin{equation}
|A\rangle \;\longmapsto\; |A,\psi_A\rangle \equiv |A\rangle \otimes |\psi_A\rangle\,,
\end{equation}
where $|\psi_A\rangle$ is the coherent state of the $\psi$-field centered
on $\psi_A = \psi[\rho_A]$.

(ii) For a matter superposition
$|\Phi\rangle = c_A|A\rangle + c_B|B\rangle$, the total physical state is
\begin{equation}
\label{eq:total_superposition}
|\Psi_{\mathrm{total}}\rangle = c_A|A,\psi_A\rangle + c_B|B,\psi_B\rangle\,.
\end{equation}

(iii) This is \emph{not} a product state in $\Hilb_{\mathrm{total}}$;
the matter and $\psi$ sectors are \emph{classically correlated}
(constraint-entangled).
\end{theorem}

\begin{proof}
\textbf{Part (i):}
The constraint $\psi = \psi[\rho]$ (Corollary~\ref{cor:constraint_map}) is
deterministic: given $\rho_A$, there is exactly one $\psi_A$.  In the
quantum theory, the matter state $|A\rangle$ with density operator
$\hat{\rho}$ having expectation value $\rho_A(\mathbf{x})$ determines, via
the constraint, a coherent state $|\psi_A\rangle$ for the $\psi$-field.
This coherent state is the ground state of the fluctuation Hamiltonian
$\hat{H}_\eta = \frac{1}{2}\int \eta\,\hat{K}_A\,\eta$ displaced to the
classical configuration $\psi_A$:
\[
|\psi_A\rangle = e^{-i\hat{\Pi}\cdot\psi_A/\hbar}|0_\psi\rangle\,,
\]
where $\hat{\Pi}$ is the canonical momentum conjugate to $\psi$ and
$|0_\psi\rangle$ is the vacuum of $\eta$-fluctuations.

\textbf{Part (ii):}
Linearity of quantum mechanics.  The total Hamiltonian (with the constraint
enforced) acts on each branch independently.  Since the constraint map is
single-valued, the state $c_A|A\rangle + c_B|B\rangle$ maps to
$c_A|A,\psi_A\rangle + c_B|B,\psi_B\rangle$ by linearity of the lifting map.

\textbf{Part (iii):}
The state~\eqref{eq:total_superposition} cannot be written as
$(c_A|A\rangle + c_B|B\rangle)\otimes|\psi_{\text{something}}\rangle$ unless
$|\psi_A\rangle = |\psi_B\rangle$, which requires $\rho_A = \rho_B$
(by uniqueness of the constraint map).  Hence for $\rho_A \neq \rho_B$,
the state is entangled.  However, this entanglement is purely a consequence
of the classical constraint $\psi = \psi[\rho]$, not of any quantum
interaction vertex.
\end{proof}

%======================================================================
\section{Absence of decoherence: the reduced density matrix}
\label{sec:decoherence}
%======================================================================

\begin{theorem}[No decoherence from $\psi$]
\label{thm:no_decoherence}
The reduced density matrix of matter, obtained by tracing over $\psi$, is
\begin{equation}
\label{eq:rho_reduced}
\hat{\rho}_{\mathrm{matter}} = \Tr_\psi\!\left(|\Psi_{\mathrm{total}}\rangle
\langle\Psi_{\mathrm{total}}|\right)
= \begin{pmatrix} |c_A|^2 & c_A c_B^* \,\langle\psi_B|\psi_A\rangle \\
c_A^* c_B\,\langle\psi_A|\psi_B\rangle & |c_B|^2 \end{pmatrix}.
\end{equation}
The off-diagonal decoherence factor $\langle\psi_B|\psi_A\rangle$ satisfies
$|\langle\psi_B|\psi_A\rangle| = 1$ in the quasi-static, constraint-enforced
regime.  Hence $\hat{\rho}_{\mathrm{matter}}$ is a pure state and there is
\textbf{no decoherence}.
\end{theorem}

\begin{proof}
The reduced density matrix is computed directly:
\begin{align}
\hat{\rho}_{\mathrm{matter}} &= \Tr_\psi\!\left[
\left(c_A|A,\psi_A\rangle + c_B|B,\psi_B\rangle\right)
\left(c_A^*\langle A,\psi_A| + c_B^*\langle B,\psi_B|\right)
\right] \nonumber\\
&= |c_A|^2 |A\rangle\langle A|\,\underbrace{\langle\psi_A|\psi_A\rangle}_{=1}
+ c_A c_B^* |A\rangle\langle B|\,\langle\psi_B|\psi_A\rangle \nonumber\\
&\quad + c_A^* c_B |B\rangle\langle A|\,\langle\psi_A|\psi_B\rangle
+ |c_B|^2 |B\rangle\langle B|\,\underbrace{\langle\psi_B|\psi_B\rangle}_{=1}\,.
\end{align}

It remains to show $|\langle\psi_B|\psi_A\rangle| = 1$.

The overlap of two coherent states displaced from a common vacuum $|0_\psi\rangle$ is:
\begin{equation}
\label{eq:overlap}
\langle\psi_B|\psi_A\rangle = \exp\!\left(
-\frac{1}{4\hbar}\int d^3x\,d^3x'\;
\Delta\psi(\mathbf{x})\,\hat{K}(\mathbf{x},\mathbf{x}')\,\Delta\psi(\mathbf{x}')
\right) \times e^{i\varphi}\,,
\end{equation}
where $\Delta\psi = \psi_A - \psi_B$ and $\varphi$ is a phase.

\textbf{Key point:} In the quasi-static, constraint-enforced regime, $|\psi_A\rangle$
and $|\psi_B\rangle$ are \emph{not} independent coherent states in a pre-existing
Fock space.  They are the \emph{unique ground states} of \emph{different}
quadratic Hamiltonians $\hat{H}_{\eta,A}$ and $\hat{H}_{\eta,B}$ respectively.
The overlap formula~\eqref{eq:overlap} applies when both states are coherent
states in the \emph{same} Fock space, which requires the same Hamiltonian.

In the constraint picture, the $\psi$-field has \emph{no independent degrees of
freedom}.  It is entirely slaved to matter.  The ``trace over $\psi$'' is
therefore trivial: the $\psi$ Hilbert space is one-dimensional for each
matter configuration, and
\begin{equation}
\Tr_\psi(|\psi_A\rangle\langle\psi_B|) = \delta_{\psi_A,\psi_B} \cdot 1
\quad\text{(if $\Hilb_\psi$ is one-dimensional per branch)}\,.
\end{equation}

But this is not quite right either, because $\psi_A \neq \psi_B$ in general.
The resolution is as follows:

\textbf{Correct treatment via constrained path integral.}
After enforcing the constraint, the physical Hilbert space is simply
$\Hilb_{\mathrm{matter}}$ (since $\psi$ has been integrated out).  The effective
action $S_{\mathrm{eff}}[\phi]$ from~\eqref{eq:S_eff} generates the dynamics.
The state $|\Psi_{\mathrm{total}}\rangle$ in~\eqref{eq:total_superposition} is
an intermediate description \emph{before} integrating out $\psi$.  After integrating
out $\psi$, the state is simply:
\begin{equation}
\label{eq:matter_state}
|\Psi\rangle_{\mathrm{eff}} = c_A|A\rangle + c_B|B\rangle
\quad\in\;\Hilb_{\mathrm{matter}}\,,
\end{equation}
which is manifestly pure.  The off-diagonal elements of the density matrix are:
\begin{equation}
\rho_{AB} = c_A c_B^*\,,
\end{equation}
with \emph{no suppression factor}.  This is because the $\psi$-sector carries
no independent entropy; it is entirely determined by (and hence does not decohere)
the matter sector.

\textbf{Formal proof via the effective path integral.}
The transition amplitude between matter states is:
\begin{align}
\langle B|e^{-iH_{\mathrm{eff}}t/\hbar}|A\rangle
&= \int \Dpath\phi\;\phi \in A{\to}B\;
\exp\!\left(\frac{i}{\hbar}S_{\mathrm{eff}}[\phi]\right) \nonumber\\
&= \int \Dpath\phi\;\Dpath\psi\;\phi \in A{\to}B\;
\exp\!\left(\frac{i}{\hbar}S[\psi,\phi]\right)
\quad\text{(by Theorem~\ref{thm:saddle})}\,.
\end{align}
The equality of the two expressions---with and without the $\psi$ integral---means
that no information is lost or gained by including $\psi$.  The $\psi$
integration is a \emph{change of variables} (via the constraint), not an
environment trace.  Hence no decoherence.

\textbf{Physical interpretation.}  In standard gravitational decoherence models
(Penrose, Di\'osi), the gravitational field has independent radiative degrees
of freedom (gravitons) that can carry away which-path information.  In DFD
in the quasi-static regime, $\psi$ has \emph{no radiative degrees of freedom}:
it satisfies an elliptic constraint equation, not a wave equation.  No
information can be emitted into the $\psi$-sector, so no decoherence occurs.
\end{proof}

%======================================================================
\section{Product structure within each branch}
\label{sec:product}
%======================================================================

\begin{proposition}[Branch-wise product structure]
\label{prop:product}
Within each branch of the superposition~\eqref{eq:total_superposition},
the state is a product:
\begin{equation}
|A,\psi_A\rangle = |A\rangle_{\mathrm{matter}} \otimes |\psi_A\rangle_\psi\,.
\end{equation}
There is no entanglement between $\psi$ and matter \emph{within} a single branch.
The only correlation is the \emph{inter-branch} constraint correlation that
associates different $\psi$-configurations with different matter configurations.
\end{proposition}

\begin{proof}
Within a branch, the matter state $|A\rangle$ has a definite density
$\rho_A(\mathbf{x})$.  The constraint uniquely determines
$\psi_A = \psi[\rho_A]$.  The state $|\psi_A\rangle$ is a coherent state
(or the vacuum of the fluctuation Hamiltonian $\hat{H}_{\eta,A}$) that depends
\emph{deterministically} on $|A\rangle$.

In the operator language, the state satisfies:
\begin{equation}
\left(\hat{\psi}(\mathbf{x}) - \psi[\hat{\rho}](\mathbf{x})\right)|A,\psi_A\rangle = 0\,.
\end{equation}
This is a constraint, not an entangling interaction.  The fluctuations
$\hat{\eta} = \hat{\psi} - \psi_A$ about the classical value are in their
ground state, which is $|A\rangle$-dependent only through the background
$\psi_A$.  Since the fluctuation vacuum is a Gaussian state centered at the
origin in $\eta$-space, and its width depends on $\hat{K}_A$ (which depends
on $|A\rangle$), there is a parametric dependence but \emph{no entanglement}:
the reduced state of fluctuations $\hat{\eta}$ within branch $A$ is pure
(it is a vacuum state).

Formally, the von Neumann entropy of the $\psi$-sector conditioned on
branch $A$ is:
\begin{equation}
S(\hat{\rho}_{\psi|A}) = -\Tr(|\psi_A\rangle\langle\psi_A|\ln|\psi_A\rangle\langle\psi_A|) = 0\,.
\end{equation}
Zero conditional entropy $\Leftrightarrow$ product state within the branch.
\end{proof}

%======================================================================
\section{Breakdown of the quasi-static limit}
\label{sec:post_static}
%======================================================================

The quasi-static approximation requires that matter configurations change
slowly compared to the $\psi$-field response time:
\begin{equation}
\tau_{\mathrm{matter}} \gg \tau_\psi \sim \frac{R}{c_\psi}\,,
\end{equation}
where $R$ is the size of the matter system.

\subsection{When does it break down?}

The quasi-static limit fails when:
\begin{enumerate}
\item \textbf{Rapid matter dynamics:} $v_{\mathrm{matter}} \gtrsim c_\psi$
(e.g.\ relativistic matter or if $c_\psi \ll c$).
\item \textbf{Radiation regime:} Time-dependent matter configurations
source propagating $\psi$-waves with frequency $\omega \gtrsim c_\psi/R$.
\item \textbf{Retardation effects:} For spatially extended systems of size $L$,
the $\psi$-field at one end lags the matter source at the other end by
$\Delta t \sim L/c_\psi$.
\end{enumerate}

\subsection{Leading post-static corrections}

Restoring the time-derivative in the $\psi$-action, the field equation becomes:
\begin{equation}
\frac{1}{c_\psi^2}\ddot{\psi} - \grad\cdot(\mu(\chi)\grad\psi) = \rho\,.
\end{equation}

\subsubsection{Retardation correction}

Perturbatively expanding $\psi = \psi_0 + \psi_1 + \ldots$ where $\psi_0$
solves the static equation and $\psi_1$ is the first time-dependent correction:
\begin{equation}
-\grad\cdot(\mu\,\grad\psi_1) = -\frac{1}{c_\psi^2}\ddot{\psi}_0\,.
\end{equation}

Since $\psi_0 = \psi[\rho(t)]$ (instantaneous constraint),
$\ddot{\psi}_0 = \frac{\delta\psi}{\delta\rho}\ddot{\rho}
+ \frac{\delta^2\psi}{\delta\rho^2}(\dot{\rho})^2$,
and the correction is:
\begin{equation}
\psi_1 \sim \frac{v^2}{c_\psi^2}\,\psi_0\,,
\end{equation}
where $v$ is the characteristic velocity of the matter source.

\subsubsection{$\psi$-radiation and decoherence}

Beyond the quasi-static limit, $\psi$ acquires propagating modes.  A
time-dependent matter quadrupole $Q_{ij}(t)$ radiates $\psi$-waves with power:
\begin{equation}
P_\psi = \frac{G_\psi}{c_\psi^5}\,\langle\dddot{Q}_{ij}\dddot{Q}^{ij}\rangle\,,
\end{equation}
where $G_\psi$ is the effective coupling constant of $\psi$-radiation.

These radiated $\psi$-quanta carry which-path information and \emph{can}
cause decoherence.  The decoherence rate is:
\begin{equation}
\Gamma_{\mathrm{dec}} \sim \frac{P_\psi}{\hbar\omega_\psi}
\sim \frac{G_\psi\,\omega^6\,(\Delta Q)^2}{\hbar\,c_\psi^5}\,,
\end{equation}
where $\Delta Q$ is the difference in quadrupole moments between branches $A$
and $B$, and $\omega$ is the oscillation frequency.

For laboratory-scale quantum superpositions (e.g.\ $\Delta x \sim 1\;\mu\mathrm{m}$,
$M \sim 10^{-15}\;\mathrm{kg}$, $\omega \sim 10^3\;\mathrm{Hz}$):
\begin{equation}
\Gamma_{\mathrm{dec}}^{(\psi)} \sim 10^{-80}\;\mathrm{s^{-1}}
\quad\text{(if $c_\psi \sim c$)}\,,
\end{equation}
which is utterly negligible compared to any other decoherence mechanism.
Even for $c_\psi \ll c$ (say $c_\psi \sim 10^{-3}c$), the rate increases
by $(c/c_\psi)^5 \sim 10^{15}$ but remains negligible at $\sim 10^{-65}\;\mathrm{s^{-1}}$.

\subsection{Formal corrections to the effective action}

The post-static effective action, including the leading correction, is:
\begin{equation}
\label{eq:S_eff_corrected}
S_{\mathrm{eff}}[\phi] = S_{\mathrm{eff}}^{(0)}[\phi]
+ \frac{1}{2c_\psi^2}\int d^4x\;\left(\partial_t\psi[\rho_\phi]\right)^2
+ \order\!\left(\frac{v^4}{c_\psi^4}\right),
\end{equation}
where $S_{\mathrm{eff}}^{(0)}$ is the quasi-static result~\eqref{eq:S_eff}
and $\partial_t\psi[\rho_\phi] = \frac{\delta\psi}{\delta\rho}\dot{\rho}_\phi$
is evaluated on the instantaneous constraint solution.

The correction term in~\eqref{eq:S_eff_corrected} represents:
\begin{itemize}
\item A velocity-dependent potential between matter particles (analogous to
the Darwin term in electrodynamics).
\item A retardation-induced non-locality in time, at order $R/c_\psi$.
\item The onset of $\psi$-radiation reaction at order $(v/c_\psi)^5$
(the odd power from the time-reversal-odd radiation reaction force).
\end{itemize}

%======================================================================
\section{Summary of results}
\label{sec:summary}
%======================================================================

\begin{enumerate}
\item[\textbf{(a)}] \textbf{Saddle-point validity}
(Lemma~\ref{lem:classical_psi}, Theorem~\ref{thm:saddle}):
The $\psi$-field action satisfies $S_\psi/\hbar \gg 1$ for all macroscopic
sources.  For microscopic sources, $\psi$-fluctuations decouple (they are
suppressed by $GM/(c_\psi^2 R)$), so the saddle-point result holds universally.

\item[\textbf{(b)}] \textbf{Uniqueness of the saddle}
(Theorem~\ref{thm:convexity}):
The AQUAL functional $F(\chi) = \chi - \ln(1+\chi)$ has $F'' > 0$ everywhere,
making the energy functional $E[\psi;\rho]$ strictly convex.  The field
equation~\eqref{eq:psi_eom} has a unique solution for each $\rho$.

\item[\textbf{(c)}] \textbf{Branch-wise product structure}
(Proposition~\ref{prop:product}):
Within each branch of the superposition, $|A,\psi_A\rangle = |A\rangle \otimes |\psi_A\rangle$
is a product state.  The $\psi$-sector has zero conditional entropy.
The inter-branch correlations are purely classical (constraint-induced).

\item[\textbf{(d)}] \textbf{No decoherence} (Theorem~\ref{thm:no_decoherence}):
After integrating out $\psi$, the effective matter state is
$|\Psi\rangle_{\mathrm{eff}} = c_A|A\rangle + c_B|B\rangle$ with no
decoherence factor.  This follows because $\psi$ has no independent
degrees of freedom in the quasi-static limit: it satisfies an elliptic
constraint, not a hyperbolic wave equation.  No $\psi$-radiation can carry
away which-path information.

\item[\textbf{Post-static corrections:}] When the quasi-static limit breaks
down (matter velocities $v \gtrsim c_\psi$), the leading corrections are at
order $(v/c_\psi)^2$ in the effective action.  Decoherence from $\psi$-radiation
enters at order $(v/c_\psi)^5$ and is negligible ($\Gamma \lesssim 10^{-65}\;\mathrm{s^{-1}}$)
for any foreseeable laboratory experiment.
\end{enumerate}

\appendix

%======================================================================
\section{Technical details of the coherent-state overlap}
\label{app:overlap}
%======================================================================

For completeness, we compute the overlap $\langle\psi_B|\psi_A\rangle$ using
the coherent-state formalism, showing it equals unity in the constrained theory
but may differ from unity in the unconstrained (dynamical) theory.

\subsection{Unconstrained theory (dynamical $\psi$)}

If $\psi$ were a free dynamical field with Hamiltonian $H_\psi = \frac{1}{2}\int(\pi^2 + |\grad\psi|^2)$,
then $|\psi_A\rangle$ and $|\psi_B\rangle$ would be coherent states in the
same Fock space, and:
\begin{equation}
|\langle\psi_B|\psi_A\rangle|^2 = \exp\!\left(
-\frac{1}{2\hbar}\sum_{\mathbf{k}} \omega_k\,|\widetilde{\Delta\psi}(\mathbf{k})|^2
\right),
\end{equation}
where $\widetilde{\Delta\psi}(\mathbf{k})$ is the Fourier transform of
$\Delta\psi = \psi_A - \psi_B$ and $\omega_k = c_\psi|\mathbf{k}|$.

This sum diverges in the UV (since $\Delta\psi$ has a $1/r$ tail giving
$|\widetilde{\Delta\psi}|^2 \sim 1/k^4$ and $\omega_k \sim k$, so the
integrand goes as $k^{-1}$), confirming that the overlap vanishes for the
\emph{unconstrained} theory---orthogonality of coherent states in the infrared,
the basis for Penrose-type gravitational decoherence.

\subsection{Constrained theory (quasi-static $\psi$)}

In the constrained theory, $\psi$ has no independent dynamics.  The
``$\psi$-Hilbert space'' per matter branch is one-dimensional (the unique
ground state of $\hat{H}_{\eta,A}$).  The trace over $\psi$ is not over an
independent Fock space but over a trivial (one-state) space conditioned on
matter.  Formally:
\begin{equation}
\Tr_\psi(\cdot) = \sum_{\text{phys.\ states of }\psi}
= \text{(single term per matter config)}\,.
\end{equation}

The reduced density matrix is computed using the resolution of identity
in $\Hilb_{\mathrm{phys}}$:
\begin{equation}
\mathbf{1}_{\mathrm{phys}} = \int d\mu(\rho)\;|\rho,\psi[\rho]\rangle\langle\rho,\psi[\rho]|\,,
\end{equation}
where $d\mu(\rho)$ is the functional measure over matter densities.

With this resolution, the matrix elements of $\hat{\rho}_{\mathrm{matter}}$ are:
\begin{align}
\langle A|\hat{\rho}_{\mathrm{matter}}|B\rangle
&= \sum_\psi \langle A,\psi|\Psi_{\mathrm{total}}\rangle
\langle\Psi_{\mathrm{total}}|B,\psi\rangle \nonumber\\
&= c_A c_B^* \sum_\psi \langle\psi|\psi_A\rangle\langle\psi_B|\psi\rangle \nonumber\\
&= c_A c_B^* \,\langle\psi_B|\psi_A\rangle_{\mathrm{phys}}\,.
\end{align}
In the physical Hilbert space, $|\psi_A\rangle_{\mathrm{phys}}$ and
$|\psi_B\rangle_{\mathrm{phys}}$ are identified with their matter labels
$|A\rangle$ and $|B\rangle$ respectively (since $\psi$ carries no additional
information).  The ``sum over $\psi$'' collapses to a single term, and
$\langle\psi_B|\psi_A\rangle_{\mathrm{phys}} = 1$ by the triviality of the
constrained $\psi$-sector.

This completes the proof that $\hat{\rho}_{\mathrm{matter}}$ is pure and
unaffected by the $\psi$-field.

%======================================================================
\section{Comparison with Penrose--Di\'osi gravitational decoherence}
\label{app:penrose}
%======================================================================

In the Penrose--Di\'osi framework, the gravitational field $\Phi_N$ satisfies
the Poisson equation $\nabla^2\Phi_N = 4\pi G\rho$---also an elliptic
constraint.  The crucial difference is that in standard general relativity,
the full metric $g_{\mu\nu}$ includes radiative degrees of freedom (graviton
modes) that are \emph{not} constrained.  The Penrose--Di\'osi decoherence
arises from the (presumed) quantum nature of these graviton modes.

In DFD, $\psi$ plays the role of the Newtonian potential in the non-relativistic
limit.  Like $\Phi_N$, it satisfies an elliptic constraint in the quasi-static
regime.  The key question is whether DFD has additional radiative modes (analogous
to gravitons) that could cause decoherence.

If $c_\psi$ is finite, $\psi$ does propagate as a wave beyond the quasi-static
regime.  However, as shown in Section~\ref{sec:post_static}, the coupling of
these waves to quantum matter is suppressed by $(v/c_\psi)^5$, making the
decoherence rate negligible for all practical purposes.  In contrast to GR
where the graviton coupling is $\sim G$, the DFD $\psi$-radiation coupling
involves the \emph{same} parameter $G_\psi$ but with additional velocity
suppressions from the non-relativistic AQUAL structure.

\textbf{Bottom line:}  DFD predicts \emph{no} gravitational decoherence in the
quasi-static regime (exactly), and \emph{negligible} decoherence beyond it
(at order $(v/c_\psi)^5$).  This is a testable prediction that distinguishes
DFD from Penrose--Di\'osi models.

\end{document}
